\ParteLibro{Comprensión de problemas y diseño de soluciones}
  {Del análisis de situaciones cotidianas a algoritmos que pueden comprobarse antes de programar.}

\chapter{Fundamentos del pensamiento computacional y análisis de problemas}
\label{cap:01-pensamiento-computacional}
\chaptermark{Pensamiento computacional}

La capacidad de resolver problemas representa uno de los pilares fundamentales de las ciencias de la computación. Antes de construir un programa informático, es necesario comprender la situación que se desea solucionar, identificar sus componentes y establecer una estrategia que permita alcanzar los resultados esperados.

Este capítulo introduce los fundamentos del pensamiento computacional como una metodología de razonamiento aplicable a diferentes situaciones de la vida cotidiana y del ámbito profesional.

\begin{resultados}
  \begin{itemize}
    \item Distinguir una necesidad cotidiana de un problema formulado para una computadora.
    \item Aplicar descomposición, reconocimiento de patrones, abstracción y generalización.
    \item Especificar objetivos, entradas, procesos, salidas y criterios de aceptación.
  \end{itemize}
\end{resultados}

\section{La resolución de problemas en la vida cotidiana}

Las personas se enfrentan constantemente a situaciones que requieren tomar decisiones y encontrar soluciones. Organizar un horario de estudios, administrar un presupuesto familiar, seleccionar una ruta de transporte o distribuir las actividades de un proyecto son ejemplos de problemas que requieren razonamiento y planificación.

Aunque estas actividades no siempre involucran el uso de computadoras, poseen características que permiten estudiarlas desde una perspectiva computacional.

\subsection{Concepto de problema}

Un problema puede definirse como una situación en la que existe una diferencia entre un estado actual y un estado deseado, cuya transformación requiere la aplicación de determinadas acciones o procedimientos. La \RefFigura{fig:elementos-problema} representa esta transición.

\FiguraLibro{02-elementos-problema}{Elementos generales de un problema}{fig:elementos-problema}{Elaboración propia.}

Desde esta perspectiva, resolver un problema implica identificar la situación inicial, establecer el objetivo y determinar las acciones necesarias para alcanzarlo.

\Needspace{11\baselineskip}
\subsection{Elementos de un problema}

\begin{tablalibro}{2}{L{0.28}L{0.72}}
\caption{Elementos de un problema}\label{tab:elementos-de-un-problema}\\
\toprule
\rowcolor{Papel}
\CabeceraTabla{Elemento} & \CabeceraTabla{Descripción} \\
\midrule
\endfirsthead
\SigueTabla{2}
\toprule
\rowcolor{Papel}
\CabeceraTabla{Elemento} & \CabeceraTabla{Descripción} \\
\midrule
\endhead
\midrule
\ContinuaTabla{2}
\endfoot
\bottomrule
\endlastfoot

Estado inicial & Situación en la que se encuentra el problema. \\
Objetivo & Resultado que se desea alcanzar. \\
Datos disponibles & Información necesaria para formular una solución. \\
Restricciones & Condiciones que deben respetarse. \\
Estrategia & Procedimientos seleccionados para resolver el problema. \\
Resultado & Producto obtenido después de aplicar la solución. \\
\end{tablalibro}


\subsection{Ejemplo práctico: administración de un presupuesto semanal}

Un estudiante dispone de S/ 100 para cubrir sus gastos semanales. Durante la semana necesita destinar dinero a transporte, alimentación y materiales educativos. Su objetivo es determinar cuánto dinero le quedará disponible después de cubrir dichos gastos. La \RefTabla{tab:presupuesto-semanal} reúne los importes del ejemplo.

\begin{tablalibro}{2}{L{0.69}R{0.31}}
\caption{Presupuesto semanal}\label{tab:presupuesto-semanal}\\
\toprule
\rowcolor{Papel}
\CabeceraTabla{Concepto} & \CabeceraTabla{Importe} \\
\midrule
\endfirsthead
\SigueTabla{2}
\toprule
\rowcolor{Papel}
\CabeceraTabla{Concepto} & \CabeceraTabla{Importe} \\
\midrule
\endhead
\midrule
\endfoot
\bottomrule
\endlastfoot

Presupuesto inicial & S/ 100 \\
Transporte & S/ 25 \\
Alimentación & S/ 40 \\
Materiales educativos & S/ 15 \\
\textbf{Total de gastos} & \textbf{S/ 80} \\
\textbf{Saldo disponible} & \textbf{S/ 20} \\
\end{tablalibro}


La solución consiste en identificar los gastos, calcular su suma y restar el resultado del presupuesto inicial.

El saldo se obtiene con la siguiente expresión.

\[ S = P - (T + A + M) \]

Las variables representan el saldo, el presupuesto inicial y los tres gastos.

\begin{itemize}
\item \(S\): saldo disponible.
\item \(P\): presupuesto inicial.
\item \(T\): gasto de transporte.
\item \(A\): gasto de alimentación.
\item \(M\): gasto de materiales.
\end{itemize}

Este ejemplo demuestra que una situación cotidiana puede transformarse en una solución organizada y expresarse mediante operaciones matemáticas. Posteriormente, este mismo procedimiento podrá representarse mediante un algoritmo e implementarse utilizando un lenguaje de programación.

\subsection{Etapas de la resolución de problemas}

\FiguraLibro{03-etapas-resolucion}{Etapas de la resolución de problemas}{fig:etapas-resolucion}{Elaboración propia.}

La aplicación sistemática de estas etapas facilita el desarrollo de soluciones organizadas. Cuando los resultados obtenidos no coinciden con los objetivos establecidos, es necesario revisar las decisiones adoptadas, identificar posibles errores y modificar la estrategia.

Por consiguiente, la resolución de problemas debe comprenderse como un proceso que puede requerir múltiples intentos, evaluaciones y mejoras.

\section{Problemas que pueden resolverse mediante computadoras}

No todos los problemas de la vida cotidiana requieren una computadora para ser resueltos. Sin embargo, cuando una situación involucra operaciones repetitivas, grandes cantidades de información o decisiones que pueden expresarse mediante condiciones precisas, el uso de un sistema computacional puede facilitar considerablemente su solución. Para determinar si resulta conveniente emplear una computadora, es necesario examinar la naturaleza del problema y la posibilidad de representar sus datos y procedimientos de manera formal.

\textcite[pp.~4–6]{farrell2013} explica que la ejecución de un programa implica realizar instrucciones y que, habitualmente, estas permiten recibir datos, procesarlos y producir resultados. También distingue los errores de sintaxis de los errores lógicos: un programa puede ejecutarse sin que su resultado sea correcto. Esta distinción permite comprender que la capacidad de procesamiento de una computadora no reemplaza el análisis previo del problema ni el diseño cuidadoso de la solución.

\subsection{Características de los problemas computacionales}

Un problema computacional es una situación que puede formularse de manera que una computadora ejecute un procedimiento para obtener una respuesta, generar un resultado o contribuir a una decisión. Para los problemas introductorios que se estudian en este libro, conviene reconocer las siguientes características:

\begin{itemize}
\item Objetivo definido. Se puede establecer qué resultado se espera obtener.
\item Datos representables. La información necesaria puede expresarse mediante números, textos, valores lógicos u otras representaciones digitales.
\item Reglas identificables. Es posible especificar operaciones, condiciones o criterios para transformar los datos.
\item Procedimiento ejecutable. Las operaciones pueden traducirse en instrucciones que una computadora sea capaz de realizar.
\item Resultado evaluable. Se dispone de criterios para comprobar si la respuesta obtenida satisface el objetivo.
\end{itemize}

Estas características constituyen una guía práctica de análisis, no una definición matemática exhaustiva de la computabilidad. Un problema puede admitir una formulación computacional y, aun así, resultar demasiado costoso de resolver exactamente con los recursos disponibles. En otros casos, la dificultad consiste en precisar los criterios o conseguir datos adecuados.

\subsection{Tipos de problemas que admiten soluciones computacionales}

Los problemas pueden clasificarse de acuerdo con la operación predominante que requiere su solución. La siguiente tabla presenta una clasificación didáctica con ejemplos de diferentes ámbitos. Las categorías no son excluyentes: una misma aplicación puede combinar varias de ellas.

\begin{tablalibro}{2}{L{0.31}L{0.69}}
\caption{Clasificación de problemas computacionales según su finalidad}\label{tab:clasificacion-de-problemas-computacionales-segun-su-finalidad}\\
\toprule
\rowcolor{Papel}
\CabeceraTabla{Tipo de problema} & \CabeceraTabla{Ejemplo de aplicación} \\
\midrule
\endfirsthead
\SigueTabla{2}
\toprule
\rowcolor{Papel}
\CabeceraTabla{Tipo de problema} & \CabeceraTabla{Ejemplo de aplicación} \\
\midrule
\endhead
\midrule
\ContinuaTabla{2}
\endfoot
\bottomrule
\multicolumn{2}{@{}p{\linewidth}@{}}{\NotaLibro{Clasificación y ejemplos de elaboración propia.}}\\
\endlastfoot

Cálculo numérico & Determinar el importe de una compra a partir de los precios y las cantidades. \\
Búsqueda y recuperación & Localizar el registro de un estudiante mediante su código. \\
Ordenamiento & Organizar productos de menor a mayor precio. \\
Clasificación y decisión & Determinar si una calificación cumple una condición de aprobación. \\
Automatización de tareas & Generar reportes periódicos a partir de registros de ventas. \\
Optimización & Seleccionar una distribución de recursos que reduzca costos bajo ciertas restricciones. \\
\end{tablalibro}


\subsection{Ejemplo práctico: control de existencias en una tienda}

Una tienda necesita identificar los productos que deben reponerse. El encargado revisa diariamente las cantidades disponibles y las compara con un nivel mínimo establecido para cada producto. Cuando la cantidad disponible es inferior al mínimo, el producto debe incorporarse a una lista de reposición.

En lugar de efectuar manualmente esta comparación para todos los productos, se propone una solución computacional que reciba los registros de existencias, compare cada cantidad con su mínimo y genere un listado con los productos que requieren atención.

Entrada: Código del producto, cantidad disponible y cantidad mínima.

Procesamiento: Comparar la cantidad disponible con la cantidad mínima de cada producto.

Salida: Listado de productos cuya cantidad disponible es inferior al mínimo.

\begin{tablalibro}{4}{L{0.39}R{0.19}R{0.17}C{0.25}}
\caption{Datos de prueba para el control de existencias}\label{tab:datos-de-prueba-para-el-control-de-existencias}\\
\toprule
\rowcolor{Papel}
\CabeceraTabla{Producto} & \CabeceraTabla{Disponible} & \CabeceraTabla{Mínimo} & \CabeceraTabla{¿Reponer?} \\
\midrule
\endfirsthead
\SigueTabla{4}
\toprule
\rowcolor{Papel}
\CabeceraTabla{Producto} & \CabeceraTabla{Disponible} & \CabeceraTabla{Mínimo} & \CabeceraTabla{¿Reponer?} \\
\midrule
\endhead
\midrule
\endfoot
\bottomrule
\endlastfoot

Cuadernos & 12 & 20 & Sí \\
Lápices & 35 & 15 & No \\
Borradores & 10 & 10 & No \\
Reglas & 4 & 8 & Sí \\
\end{tablalibro}


El resultado esperado es una lista que contiene cuadernos y reglas. Obsérvese que los borradores no se incluyen, porque la regla establecida exige que la cantidad disponible sea estrictamente inferior al mínimo. Si el encargado quisiera reponer también los productos que tienen exactamente la cantidad mínima, sería necesario modificar la condición. Esta diferencia muestra cómo una pequeña variación en la formulación del problema puede cambiar los resultados del programa.

\subsection{Alcances y limitaciones de las soluciones computacionales}

Una computadora puede ejecutar operaciones con rapidez, almacenar información y aplicar criterios de manera consistente. No obstante, un sistema basado en reglas depende de la calidad de sus datos y de la corrección de los procedimientos definidos. Por ejemplo, si el inventario registra una cantidad equivocada, la lista de reposición podría ser incorrecta aunque el programa realice correctamente todas las comparaciones.

Asimismo, existen situaciones en las que no basta con aplicar reglas explícitas. La identificación de objetos en imágenes, por ejemplo, puede requerir modelos que aprendan a partir de datos. Sus resultados no necesariamente son exactos y deben evaluarse antes de emplearlos en decisiones relevantes. Este ejemplo permite distinguir los procedimientos deterministas sencillos, que serán el centro de los primeros capítulos, de otras estrategias computacionales más avanzadas.

Por tanto, antes de decidir que una computadora debe resolver un problema, conviene preguntarse si el objetivo está suficientemente definido, si los datos están disponibles, si el procedimiento es viable y cómo se verificará el resultado. Estas preguntas preparan al estudiante para comprender el pensamiento computacional, tema del siguiente apartado.

\section{Concepto e importancia del pensamiento computacional}

El pensamiento computacional emplea conceptos de la informática para formular problemas y organizar soluciones. \textcite{wing2006} destaca dos recursos para tratar tareas complejas: descomponerlas y elegir una representación que conserve lo relevante. Este capítulo combina esas prácticas con el reconocimiento de patrones y la generalización. La selección de estrategias depende de qué se necesita resolver y cómo se puede comprobar el resultado.

La planificación de la lógica constituye una etapa anterior a la codificación. \textcite[pp.~8–10]{farrell2013} explica que el desarrollo de programas exige comprender las necesidades del usuario, determinar el resultado esperado y organizar las operaciones que conducen a dicho resultado. Esta perspectiva ofrece un punto de partida para relacionar el pensamiento computacional con la práctica de la programación.

\subsection{Componentes del pensamiento computacional}

Para los fines didácticos de este libro, se distinguen cinco estrategias complementarias. No representan una secuencia rígida: durante el análisis es frecuente regresar a una estrategia anterior cuando aparece información nueva o se descubre una condición que no había sido considerada.

\begin{tablalibro}{3}{L{0.28}L{0.34}L{0.38}}
\caption{Estrategias del pensamiento computacional}\label{tab:estrategias-del-pensamiento-computacional}\\
\toprule
\rowcolor{Papel}
\CabeceraTabla{Estrategia} & \CabeceraTabla{Pregunta orientadora} & \CabeceraTabla{Ejemplo} \\
\midrule
\endfirsthead
\SigueTabla{3}
\toprule
\rowcolor{Papel}
\CabeceraTabla{Estrategia} & \CabeceraTabla{Pregunta orientadora} & \CabeceraTabla{Ejemplo} \\
\midrule
\endhead
\midrule
\endfoot
\bottomrule
\multicolumn{3}{@{}p{\linewidth}@{}}{\NotaLibro{Elaboración propia.}}\\
\endlastfoot

Descomposición & ¿En qué partes puedo dividir el problema? & Separar registro, cálculo y reporte de notas. \\
Reconocimiento de patrones & ¿Qué se repite o se parece? & Identificar cálculos iguales en distintos cursos. \\
Abstracción & ¿Qué información necesito conservar? & Considerar las calificaciones y sus pesos. \\
Generalización & ¿Cómo reutilizar la solución? & Aplicar el cálculo a distintos estudiantes. \\
Diseño algorítmico & ¿Qué pasos permiten obtener el resultado? & Validar datos, calcular y comunicar el promedio. \\
\end{tablalibro}


\subsection{Importancia en la formación académica y profesional}

El pensamiento computacional ayuda a organizar problemas complejos, explicar decisiones y establecer criterios de comprobación. En el ámbito educativo permite estructurar investigaciones, analizar información y construir modelos sencillos. En ingeniería facilita la especificación de procesos y la identificación de restricciones. En administración contribuye a organizar tareas repetitivas y a detectar inconsistencias en registros.

Su utilidad no implica que todo problema deba automatizarse. Una solución puede ser técnicamente posible y, aun así, resultar inadecuada por su costo, por la falta de datos confiables o por sus consecuencias para las personas. El análisis debe considerar tanto la viabilidad técnica como las necesidades de quienes utilizarán el resultado.

\subsection{Ejemplo: organización de una jornada de estudio}

Un estudiante dispone de 180 minutos para preparar tres asignaturas. Decide reservar 15 minutos de descanso y distribuir los 165 minutos restantes según la dificultad de cada asignatura. Si establece pesos de 3, 2 y 1, puede asignar 82,5; 55 y 27,5 minutos, respectivamente. La solución identifica recursos, restricciones y una regla de distribución. Si el tiempo disponible cambia, la misma regla puede aplicarse con nuevos datos.

Antes de automatizar esta distribución, conviene verificar que los pesos sean positivos, que el descanso no exceda el tiempo total y que el reparto sea compatible con las necesidades reales del estudiante. El ejemplo muestra que razonar computacionalmente incluye formular condiciones y revisar resultados, además de efectuar cálculos.

\section{Descomposición de problemas en partes más pequeñas}

La descomposición consiste en dividir un problema complejo en subproblemas cuya comprensión y tratamiento resulten más manejables. Cada parte debe tener una finalidad identificable y contribuir al objetivo general. Esta estrategia reduce la cantidad de detalles que deben considerarse simultáneamente y facilita localizar las dificultades que requieren mayor atención.

En el diseño de programas, \textcite[pp.~48–50]{farrell2013} relaciona la división de tareas con la modularización, la abstracción y la reutilización. Aunque en este capítulo todavía no se implementan módulos de código, la descomposición permite preparar una estructura que más adelante podrá traducirse en componentes de un programa.

\subsection{Procedimiento de descomposición}

\begin{itemize}
\item Precisar el objetivo general y delimitar qué queda fuera del problema.
\item Identificar las tareas principales y el resultado de cada una.
\item Dividir las tareas complejas hasta obtener actividades comprensibles.
\item Establecer las dependencias entre las partes y la información que intercambian.
\item Revisar que la integración de las partes satisfaga el objetivo original.
\end{itemize}

\subsection{Ejemplo: sistema de préstamo de libros}

Una biblioteca escolar necesita registrar préstamos y devoluciones. Intentar resolver todas las operaciones como una sola tarea dificulta identificar las reglas. Es preferible distinguir la identificación del lector, la consulta de disponibilidad, el registro del préstamo y la actualización de la devolución.

\begin{tablalibro}{3}{L{0.28}L{0.34}L{0.38}}
\caption{Descomposición de un sistema de préstamos}\label{tab:descomposicion-de-un-sistema-de-prestamos}\\
\toprule
\rowcolor{Papel}
\CabeceraTabla{Subproblema} & \CabeceraTabla{Información requerida} & \CabeceraTabla{Resultado} \\
\midrule
\endfirsthead
\SigueTabla{3}
\toprule
\rowcolor{Papel}
\CabeceraTabla{Subproblema} & \CabeceraTabla{Información requerida} & \CabeceraTabla{Resultado} \\
\midrule
\endhead
\midrule
\endfoot
\bottomrule
\multicolumn{3}{@{}p{\linewidth}@{}}{\NotaLibro{Elaboración propia.}}\\
\endlastfoot

Identificar lector & Código del lector & Lector reconocido o aviso de error \\
Consultar ejemplar & Código del ejemplar & Disponibilidad confirmada \\
Registrar préstamo & Lector, ejemplar y fecha & Préstamo registrado \\
Registrar devolución & Identificador del préstamo & Préstamo cerrado y ejemplar disponible \\
\end{tablalibro}


Las partes no son necesariamente independientes. El registro del préstamo requiere que el lector esté identificado y que el ejemplar se encuentre disponible. Si una consulta falla, no se debe modificar el estado del ejemplar. Definir estas dependencias evita que cada componente funcione correctamente por separado, pero produzca un resultado incorrecto al integrarse.

\subsection{Criterios de una buena descomposición}

Una división adecuada asigna a cada parte una responsabilidad clara, evita duplicar tareas y permite verificar resultados parciales. Una división demasiado extensa puede ocultar dificultades; una información excesivamente detallada puede generar complejidad innecesaria. El nivel de detalle debe corresponder al propósito del análisis y a la experiencia de quienes participan en él.

\section{Reconocimiento de patrones y semejanzas}

El reconocimiento de patrones consiste en identificar regularidades, relaciones o estructuras semejantes entre situaciones. En la resolución de problemas, permite aprovechar conocimientos previos y evitar diseñar desde el inicio procedimientos que ya se han aplicado con éxito en circunstancias equivalentes.

Un patrón no siempre es una repetición exacta. Dos problemas pueden utilizar datos diferentes y compartir una misma estructura lógica. Por ejemplo, calcular el total de una compra y calcular el costo de materiales de un proyecto requieren multiplicar cantidades por precios unitarios y sumar los importes obtenidos.

\Needspace{11\baselineskip}
\subsection{Patrones frecuentes en problemas computacionales}

\begin{tablalibro}{3}{L{0.28}L{0.34}L{0.38}}
\caption{Patrones de resolución y ejemplos}\label{tab:patrones-de-resolucion-y-ejemplos}\\
\toprule
\rowcolor{Papel}
\CabeceraTabla{Patrón} & \CabeceraTabla{Operación común} & \CabeceraTabla{Aplicación} \\
\midrule
\endfirsthead
\SigueTabla{3}
\toprule
\rowcolor{Papel}
\CabeceraTabla{Patrón} & \CabeceraTabla{Operación común} & \CabeceraTabla{Aplicación} \\
\midrule
\endhead
\midrule
\endfoot
\bottomrule
\endlastfoot

Acumulación & Sumar valores sucesivos & Total de gastos \\
Conteo & Incrementar un contador & Número de asistentes \\
Comparación & Evaluar una condición & Productos bajo el mínimo \\
Búsqueda & Localizar un elemento & Encontrar un código de matrícula \\
Clasificación & Asignar una categoría & Agrupar solicitudes por estado \\
\end{tablalibro}


\subsection{Ejemplo: análisis de asistencia}

Un docente registra cinco sesiones con los estados presente, ausente, presente, presente y ausente. Para obtener la cantidad de asistencias, se reconoce una operación que debe repetirse para cada registro: examinar el estado e incrementar el conteo cuando corresponda. El resultado es tres asistencias de cinco sesiones, equivalente al 60 \% si todas las sesiones tienen el mismo valor.

La misma estructura permite contar entregas de tareas o productos defectuosos, siempre que se cambie el criterio de selección. Sin embargo, si una sesión justificada debe excluirse del denominador, la regla de cálculo ya no es idéntica. Reconocer semejanzas exige también identificar diferencias que afectan la solución.

\subsection{Patrones y evidencia}

En problemas basados en datos, una regularidad observada puede orientar una hipótesis, pero no demuestra que se mantenga en todos los casos ni establece una relación causal. Por ello, los patrones deben contrastarse con ejemplos nuevos, excepciones y datos que no se utilizaron para formular la regla. Esta precaución será especialmente relevante en aplicaciones de aprendizaje automático.

\section{Abstracción e identificación de información relevante}

La abstracción permite representar un problema conservando las características que son necesarias para alcanzar un objetivo y dejando de lado detalles que no afectan la solución. \textcite[p.~49]{farrell2013} presenta la abstracción como un recurso que facilita comprender tareas complejas al concentrarse en sus propiedades esenciales.

La relevancia de un dato depende del propósito. Un dato puede ser indispensable en una aplicación y prescindible en otra. Por tanto, abstraer no significa eliminar información arbitrariamente, sino justificar qué datos se conservan, qué relaciones se representan y qué detalles se omiten.

\subsection{Ejemplo: cálculo de una tarifa de transporte}

Supóngase una tarifa didáctica compuesta por S/ 4 de cargo fijo y S/ 1,50 por kilómetro recorrido. Para calcular el importe se necesita la distancia y los valores de la tarifa. El color del vehículo o el nombre del conductor no intervienen en el cálculo. En cambio, si se incorporan recargos nocturnos, la hora del viaje pasa a ser relevante.

\begin{tablalibro}{3}{L{0.28}C{0.18}L{0.54}}
\caption{Selección de información para calcular una tarifa}\label{tab:seleccion-de-informacion-para-calcular-una-tarifa}\\
\toprule
\rowcolor{Papel}
\CabeceraTabla{Dato} & \CabeceraTabla{¿Relevante?} & \CabeceraTabla{Justificación} \\
\midrule
\endfirsthead
\SigueTabla{3}
\toprule
\rowcolor{Papel}
\CabeceraTabla{Dato} & \CabeceraTabla{¿Relevante?} & \CabeceraTabla{Justificación} \\
\midrule
\endhead
\midrule
\endfoot
\bottomrule
\endlastfoot

Distancia recorrida & Sí & Determina el cargo variable. \\
Cargo fijo & Sí & Forma parte del importe. \\
Precio por kilómetro & Sí & Determina el costo por distancia. \\
Color del vehículo & No & No modifica la tarifa definida. \\
Hora del viaje & Depende & Es necesaria si existen recargos horarios. \\
\end{tablalibro}


Con una distancia de 6 km, el importe del ejemplo es S/ 4 + S/ 1,50 × 6 = S/ 13. Esta representación supone que no existen peajes, descuentos ni otras condiciones. Si tales condiciones forman parte del problema real, deberán añadirse al modelo antes de utilizarlo.

\subsection{Niveles de abstracción}

En un nivel conceptual se describen objetivos y entidades; en un nivel lógico se identifican datos, relaciones y reglas; en un nivel de implementación se establecen instrucciones y estructuras concretas. En el presente capítulo se trabaja principalmente en los dos primeros niveles. Esta separación permite discutir el problema sin depender todavía de la sintaxis de un lenguaje.

Una abstracción insuficiente conserva detalles que dificultan comprender el problema. Una abstracción excesiva elimina condiciones necesarias y puede generar resultados incorrectos. Para evaluar una representación, debe comprobarse si permite responder las preguntas planteadas sin introducir supuestos injustificados.

\section{Generalización de soluciones para diferentes situaciones}

La generalización consiste en transformar una solución particular en un procedimiento aplicable a una familia de problemas que comparten características esenciales. Se apoya en el reconocimiento de patrones y en la abstracción: primero se identifica aquello que permanece constante y después se expresa como una regla que admite diferentes valores de entrada.

\subsection{De una solución particular a una regla general}

Considérese una compra de tres cuadernos a S/ 5 cada uno. El importe es S/ 15. Si se desea resolver compras de otros productos, no conviene conservar los valores 3 y 5 como parte fija del procedimiento. Se sustituyen por datos variables: cantidad y precio unitario. Así, el importe se obtiene multiplicando ambos valores.

Para una compra con varios productos, la regla se amplía: calcular el importe de cada producto y sumar los importes. El número de productos deja de estar fijado de antemano. El procedimiento puede utilizarse para una compra de uno, dos o más artículos, bajo las mismas condiciones de cálculo.

\begin{tablalibro}{4}{L{0.43}R{0.16}R{0.21}R{0.20}}
\caption{Aplicación de una regla general de cálculo}\label{tab:aplicacion-de-una-regla-general-de-calculo}\\
\toprule
\rowcolor{Papel}
\CabeceraTabla{Situación} & \CabeceraTabla{Cantidad} & \CabeceraTabla{Precio unitario} & \CabeceraTabla{Importe} \\
\midrule
\endfirsthead
\SigueTabla{4}
\toprule
\rowcolor{Papel}
\CabeceraTabla{Situación} & \CabeceraTabla{Cantidad} & \CabeceraTabla{Precio unitario} & \CabeceraTabla{Importe} \\
\midrule
\endhead
\midrule
\endfoot
\bottomrule
\endlastfoot

Cuadernos & 3 & S/ 5 & S/ 15 \\
Lápices & 8 & S/ 1,50 & S/ 12 \\
Carpetas & 2 & S/ 9 & S/ 18 \\
\end{tablalibro}


\subsection{Condiciones y límites de la generalización}

Generalizar no significa que una solución sea válida para cualquier problema. La regla anterior presupone cantidades no negativas y precios unitarios válidos. Si existen descuentos por volumen, impuestos diferenciados o devoluciones, es necesario ampliar las condiciones. Una solución reutilizable debe indicar su ámbito de aplicación y las situaciones que no contempla.

Para comprobar una generalización conviene ensayar casos diferentes: cantidades pequeñas, grandes, iguales a cero y datos inválidos. La finalidad es determinar si la regla conserva su significado al cambiar los valores y si responde de manera coherente a las excepciones. En los capítulos posteriores, esta práctica se relacionará con la parametrización, las funciones y las pruebas de programas.

\section{Formulación de problemas computacionales: objetivos, entradas, procesos y salidas}

Formular un problema computacional significa expresar de manera clara y verificable una necesidad que se desea atender mediante un procedimiento computacional. Una formulación adecuada establece qué debe conseguirse, qué datos se utilizarán, qué transformaciones se realizarán y cómo se reconocerá un resultado satisfactorio.

\textcite[pp.~8–10]{farrell2013} señala que, al planificar un programa, es útil comenzar por determinar la salida deseada y, a partir de ella, establecer las entradas y el procesamiento necesarios. Este enfoque evita recopilar datos sin una finalidad definida y facilita la planificación de la lógica.

\Needspace{11\baselineskip}
\subsection{Objetivos, entradas, procesos y salidas}

\begin{tablalibro}{3}{L{0.22}L{0.40}L{0.38}}
\caption{Elementos de la formulación computacional}\label{tab:elementos-de-la-formulacion-computacional}\\
\toprule
\rowcolor{Papel}
\CabeceraTabla{Elemento} & \CabeceraTabla{Descripción} & \CabeceraTabla{Pregunta orientadora} \\
\midrule
\endfirsthead
\SigueTabla{3}
\toprule
\rowcolor{Papel}
\CabeceraTabla{Elemento} & \CabeceraTabla{Descripción} & \CabeceraTabla{Pregunta orientadora} \\
\midrule
\endhead
\midrule
\endfoot
\bottomrule
\endlastfoot

Objetivo & Resultado o necesidad que se busca satisfacer. & ¿Qué se desea conseguir? \\
Entradas & Datos necesarios para realizar el procedimiento. & ¿Qué información se requiere? \\
Procesos & Operaciones y condiciones que transforman los datos. & ¿Qué se debe hacer? \\
Salidas & Información o resultados producidos. & ¿Qué debe entregar la solución? \\
Restricciones & Reglas y límites que deben respetarse. & ¿Qué condiciones se deben cumplir? \\
Criterios de aceptación & Condiciones para comprobar el resultado. & ¿Cómo sabemos que funciona? \\
\end{tablalibro}


\subsection{Ejemplo: cálculo del promedio ponderado}

Una institución ficticia requiere calcular una nota final a partir de tres evaluaciones con pesos de 30 \%, 30 \% y 40 \%. Las calificaciones se expresan en una escala de 0 a 20. Para este ejemplo se adopta una nota mínima aprobatoria de 11, sin redondeo previo a la decisión; estas reglas son supuestos didácticos y no corresponden a un reglamento institucional específico.

\begin{tablalibro}{3}{L{0.30}L{0.45}L{0.25}}
\caption{Modelo de entrada, procesamiento y salida}\label{tab:modelo-de-entrada-procesamiento-y-salida}\\
\toprule
\rowcolor{Papel}
\CabeceraTabla{Entrada} & \CabeceraTabla{Procesamiento} & \CabeceraTabla{Salida} \\
\midrule
\endfirsthead
\SigueTabla{3}
\toprule
\rowcolor{Papel}
\CabeceraTabla{Entrada} & \CabeceraTabla{Procesamiento} & \CabeceraTabla{Salida} \\
\midrule
\endhead
\midrule
\ContinuaTabla{3}
\endfoot
\bottomrule
\multicolumn{3}{@{}p{\linewidth}@{}}{\NotaLibro{Elaboración propia.}}\\
\endlastfoot

Notas: 14, 12 y 16; pesos: 0,30; 0,30; 0,40. & Validar notas y pesos; calcular 14 × 0,30 + 12 × 0,30 + 16 × 0,40; comparar con 11. & Nota final: 14,2. Condición: aprobado. \\
\end{tablalibro}


La formulación también debe indicar qué ocurre si una nota falta o está fuera del intervalo permitido. En este caso, se establece que el sistema informe el error y no emita una condición académica hasta que los datos sean válidos. Esta decisión evita confundir una entrada incorrecta con un resultado de desaprobación.

\subsection{Comprobación preliminar de la formulación}

Antes de diseñar el algoritmo conviene revisar que los objetivos no sean ambiguos, que las entradas estén disponibles, que las reglas sean consistentes y que las salidas puedan verificarse mediante ejemplos. Una especificación que dice únicamente «calcular notas» es insuficiente: debe indicar la escala, los pesos, la condición de aprobación y el tratamiento de los datos inválidos.

\begin{tablalibro}{3}{L{0.28}L{0.35}L{0.37}}
\caption{Casos de comprobación del promedio ponderado}\label{tab:casos-de-comprobacion-del-promedio-ponderado}\\
\toprule
\rowcolor{Papel}
\CabeceraTabla{Notas} & \CabeceraTabla{Resultado esperado} & \CabeceraTabla{Finalidad} \\
\midrule
\endfirsthead
\SigueTabla{3}
\toprule
\rowcolor{Papel}
\CabeceraTabla{Notas} & \CabeceraTabla{Resultado esperado} & \CabeceraTabla{Finalidad} \\
\midrule
\endhead
\midrule
\endfoot
\bottomrule
\endlastfoot

14, 12, 16 & 14,2; aprobado & Caso habitual \\
11, 11, 11 & 11; aprobado & Límite de aprobación \\
10, 10, 10 & 10; desaprobado & Caso inferior al límite \\
21, 12, 16 & Mensaje de error & Entrada fuera de rango \\
\end{tablalibro}


Estos resultados permiten verificar la coherencia del planteamiento incluso antes de escribir un programa. En el \RefCapitulo{cap:02-diseno-algoritmico} se utilizarán especificaciones de este tipo para construir algoritmos, pseudocódigo y diagramas de flujo.

\SeccionLibro[Caso práctico]{Caso práctico. Análisis de una situación cotidiana}

\subsection{Control del presupuesto de materiales de un aula}

Un aula dispone de S/ 100,00 para adquirir materiales destinados a una exposición. El responsable necesita conocer el gasto total, el saldo y la viabilidad de la compra. Cada material tiene un nombre, una cantidad entera positiva y un precio unitario no negativo con dos decimales. Una lista vacía representa una compra de importe cero. El caso utiliza los mismos datos de referencia en los cuatro capítulos.

\begin{tablalibro}{4}{L{0.37}R{0.16}R{0.25}R{0.22}}
\caption{Datos iniciales del caso práctico}\label{tab:datos-iniciales-del-caso-practico}\\
\toprule
\rowcolor{Papel}
\CabeceraTabla{Material} & \CabeceraTabla{Cantidad} & \CabeceraTabla{Precio unitario} & \CabeceraTabla{Subtotal} \\
\midrule
\endfirsthead
\SigueTabla{4}
\toprule
\rowcolor{Papel}
\CabeceraTabla{Material} & \CabeceraTabla{Cantidad} & \CabeceraTabla{Precio unitario} & \CabeceraTabla{Subtotal} \\
\midrule
\endhead
\midrule
\endfoot
\bottomrule
\endlastfoot

Cuadernos & 5 & S/ 4,50 & S/ 22,50 \\
Lápices & 10 & S/ 1,20 & S/ 12,00 \\
Cartulinas & 4 & S/ 3,00 & S/ 12,00 \\
\textbf{Total} & & & \textbf{S/ 46,50} \\
\end{tablalibro}


\subsection{Aplicación de las estrategias del capítulo}

\begin{itemize}
\item Descomposición: separar el registro de materiales, la validación, el cálculo de subtotales, la suma y la comparación con el presupuesto.
\item Reconocimiento de patrones: observar que cada subtotal resulta de multiplicar la cantidad por el precio unitario.
\item Abstracción: conservar nombre, cantidad, precio y presupuesto; omitir características que no afectan el costo, como el color de los materiales.
\item Generalización: permitir una lista de cualquier longitud y un presupuesto variable, bajo las mismas reglas.
\item Formulación: definir las entradas, las operaciones, el saldo y el mensaje que informa si la compra es viable.
\end{itemize}

El total es S/ 46,50 y el saldo es S/ 53,50. Una compra es viable cuando el saldo es cero o positivo. Si el saldo es negativo, su magnitud indica cuánto dinero falta. Un dato inválido genera un mensaje de error antes de emitir un resultado de compra.

\begin{tablalibro}{3}{R{0.25}R{0.25}L{0.50}}
\caption{Comprobación de situaciones alternativas}\label{tab:comprobacion-de-situaciones-alternativas}\\
\toprule
\rowcolor{Papel}
\CabeceraTabla{Presupuesto} & \CabeceraTabla{Gasto total} & \CabeceraTabla{Resultado esperado} \\
\midrule
\endfirsthead
\SigueTabla{3}
\toprule
\rowcolor{Papel}
\CabeceraTabla{Presupuesto} & \CabeceraTabla{Gasto total} & \CabeceraTabla{Resultado esperado} \\
\midrule
\endhead
\midrule
\ContinuaTabla{3}
\endfoot
\bottomrule
\multicolumn{3}{@{}p{\linewidth}@{}}{\NotaLibro{Elaboración propia.}}\\
\endlastfoot

S/ 150 & S/ 95 & Viable; saldo S/ 55 \\
S/ 95 & S/ 95 & Viable; saldo S/ 0 \\
S/ 80 & S/ 95 & No viable; faltan S/ 15 \\
S/ 150 & Precio negativo & Error; no calcular \\
\end{tablalibro}


El caso integra los contenidos del capítulo sin exigir conocimientos de programación. El producto de esta etapa es una especificación que puede servir como base para diseñar un algoritmo en el capítulo siguiente.

\SeccionLibro{Cierre del capítulo}

\subsection{Resumen}

La resolución de problemas comienza con la comprensión de una necesidad y la definición de un resultado esperado. El pensamiento computacional aporta estrategias para organizar este proceso. La descomposición divide tareas complejas; el reconocimiento de patrones identifica estructuras comunes; la abstracción selecciona información relevante; y la generalización facilita reutilizar procedimientos. Finalmente, la formulación computacional integra objetivos, entradas, procesos, salidas, restricciones y criterios de aceptación. Estas actividades preparan el diseño de soluciones algorítmicas y reducen la posibilidad de implementar programas que no respondan a las necesidades del usuario.

\Needspace{11\baselineskip}
\subsection{Términos clave}

\begin{tablalibro}{2}{L{0.31}L{0.69}}
\caption{Glosario del capítulo}\label{tab:glosario-del-capitulo-1}\\
\toprule
\rowcolor{Papel}
\CabeceraTabla{Término} & \CabeceraTabla{Significado} \\
\midrule
\endfirsthead
\SigueTabla{2}
\toprule
\rowcolor{Papel}
\CabeceraTabla{Término} & \CabeceraTabla{Significado} \\
\midrule
\endhead
\midrule
\endfoot
\bottomrule
\multicolumn{2}{@{}p{\linewidth}@{}}{\NotaLibro{Elaboración propia.}}\\
\endlastfoot

Problema computacional & Necesidad expresable mediante datos y procedimientos computacionales. \\
Pensamiento computacional & Razonamiento orientado a formular y organizar soluciones computacionales. \\
Descomposición & División de un problema en partes manejables. \\
Patrón & Regularidad o estructura común entre situaciones. \\
Abstracción & Selección de características relevantes para un objetivo. \\
Generalización & Extensión de una solución a casos semejantes. \\
Entrada & Dato requerido para efectuar un proceso. \\
Proceso & Transformación u operación sobre los datos. \\
Salida & Resultado producido por el procedimiento. \\
Restricción & Condición que limita la solución. \\
\end{tablalibro}


\subsection{Preguntas de comprensión}

\begin{enumerate}
\item ¿En qué se diferencia resolver un problema de escribir un programa?
\item ¿Qué características permiten formular una situación como problema computacional?
\item ¿Por qué el pensamiento computacional no se limita al uso de computadoras?
\item ¿Cómo contribuye la descomposición a comprender un problema complejo?
\item ¿Qué diferencia existe entre reconocer un patrón y generalizar una solución?
\item ¿Por qué un dato puede ser relevante para un problema e irrelevante para otro?
\item ¿Qué riesgos presenta una abstracción excesiva?
\item ¿Por qué resulta útil comenzar la formulación por el resultado esperado?
\item ¿Qué función cumplen las restricciones y los criterios de aceptación?
\item ¿Cómo se puede comprobar una formulación antes de implementar un programa?
\end{enumerate}

\subsection{Ejercicios de análisis y abstracción}

\begin{enumerate}
\item Una cafetería necesita calcular el costo de pedidos. Identifique el objetivo, los datos relevantes y tres datos prescindibles. Establezca cómo tratar cantidades negativas.
\item Una biblioteca necesita contabilizar préstamos vencidos. Descomponga el problema y establezca qué fecha se utilizará para determinar el vencimiento.
\item Un docente desea calcular el porcentaje de asistencia. Identifique el patrón de conteo y proponga una regla que considere las ausencias justificadas.
\item Una familia distribuye un presupuesto mensual entre alimentación, transporte y vivienda. Generalice el cálculo del saldo para cualquier número de categorías.
\item Una tienda compara existencias con niveles mínimos. Formule las entradas, el proceso y las salidas; incluya ejemplos de igualdad y de datos inválidos.
\end{enumerate}

\subsection{Desafíos de resolución de problemas}

\textbf{Desafío 1.} Organización de horarios. Diseñe la formulación de un sistema que distribuya actividades académicas en intervalos disponibles sin superposiciones. Identifique restricciones, subproblemas y al menos tres casos de comprobación. No es necesario escribir código.

\textbf{Desafío 2.} Registro de reciclaje. Una institución desea resumir la cantidad de residuos reciclables recolectados por semana y por categoría. Proponga una representación de datos, reglas de validación y una solución generalizable a distintos periodos. Explique qué datos dejaría fuera y por qué.

\textbf{Desafío 3.} Integración. Seleccione una necesidad de su entorno, describa el estado inicial y el objetivo, aplique las cinco estrategias de análisis y elabore una tabla de entradas, procesos y salidas. Presente tres casos válidos y uno inválido con sus resultados esperados. La evaluación considerará claridad, coherencia, tratamiento de restricciones y posibilidad de comprobar la solución.
